\documentclass[11pt,a4paper]{article}
\usepackage[margin=1in]{geometry}
\usepackage{amsmath, mathtools, amssymb, amsthm}
\usepackage{graphicx}
\usepackage{booktabs}
\usepackage{xcolor}
\usepackage[colorlinks=true, linkcolor=blue, citecolor=blue, urlcolor=blue]{hyperref}

\theoremstyle{plain}
\newtheorem{theorem}{Theorem}
\newtheorem{proposition}[theorem]{Proposition}
\theoremstyle{definition}
\newtheorem{definition}[theorem]{Definition}
\newtheorem{remark}[theorem]{Remark}

\newcommand{\todo}[1]{\vspace{0.3em}\noindent\fbox{\parbox{0.95\textwidth}{\textit{\textbf{Draft note:} #1}}}\vspace{0.3em}}
\newcommand{\SUE}{\mathrm{SUE}}
\newcommand{\AR}{\mathrm{AR}}
\newcommand{\IR}{\mathrm{IR}}
\newcommand{\IC}{\mathrm{IC}}

\title{\textbf{A Disciplined Pipeline for Weak Cross-Sectional Equity Signals:\\
Point-in-Time PEAD, Same-Industry Lead--Lag, Utility-Weighted Gating,\\
Risk-Parity Combination, and Regime-Aware Sizing}}
\author{Stefanos Drakos\\
\small AGEL AI I.K.E., Rhodes, Greece\\
\small ORCID: 0000-0001-7417-2444}
\date{\today}

\begin{document}
\maketitle

\begin{abstract}
We present a transparent, reproducible pipeline for constructing and \emph{honestly evaluating}
weak cross-sectional equity signals from free, point-in-time data. The pipeline has five stages:
(i) a time-series post-earnings-announcement-drift (PEAD) signal built only from reported earnings
via SEC EDGAR, point-in-time by construction; (ii) a same-industry lead--lag signal that routes a
firm's earnings surprise to its industry peers; (iii) a utility-weighted gate that admits a signal
only if its long--short portfolio carries economically significant, time-stable edge; (iv) a
risk-parity (equal-risk-contribution) combination of the surviving signals; and (v) a
regime-aware sizing layer that scales exposure by the prevailing volatility state. Two
methodological points are central. First, the same-industry lead--lag effect is \emph{confounded}
with own-firm PEAD through earnings-season clustering; we separate them with a daily
Fama--MacBeth cross-sectional regression rather than ad-hoc exclusion. Second, we treat
\emph{durability}---the decay of edge across calendar time---as the headline diagnostic, not the
in-sample information ratio. On real US equity data
[\emph{universe, sample, headline} $\IR$/$t_{\mathrm{NW}}$/\emph{half-life}/\emph{durability to be inserted}], the
combined signal achieves [\ldots]. We argue that the contribution is the evaluation protocol
itself: a null result with correct Newey--West inference and a reported decay profile is a valid,
publishable outcome. These results operationalize the anomalies of Bernard and Thomas (1989) and
Cohen and Frazzini (2008) under the overfitting controls of Harvey, Liu, and Zhu (2016) and Bailey
and L\'opez de Prado (2014).
\end{abstract}

\section{Introduction}
\label{sec:intro}

\paragraph{Problem.} Cross-sectional equity signals derived from freely available daily data are
individually weak (information coefficients of order $0.02$--$0.05$) and increasingly fragile: the
classical PEAD drift of roughly two percent over sixty trading days documented by Bernard and
Thomas (1989) has attenuated materially in modern samples. In this regime, model \emph{capacity}
is a liability: flexible learners fit noise and destroy otherwise robust low-capacity factors.

\paragraph{Prior work.} PEAD originates with Ball and Brown (1968) and is named and characterized
by Bernard and Thomas (1989), using a time-series (seasonal-random-walk) earnings expectation;
Foster, Olsen, and Shevlin (1984) establish the standardized-unexpected-earnings construction.
Cross-firm information diffusion---return predictability across economically linked firms---is
documented by Cohen and Frazzini (2008) and, at the industry level, by Menzly and Ozbas (2010).
Volatility-managed sizing is analyzed by Moreira and Muir (2017). On the evaluation side, Harvey,
Liu, and Zhu (2016) quantify the multiple-testing problem in factor research, and Bailey and
L\'opez de Prado (2014) introduce the Deflated Sharpe Ratio; combinatorial purged cross-validation
and the probability of backtest overfitting are developed in L\'opez de Prado (2018).

\paragraph{Contribution.} We do not claim a new anomaly. We contribute (1) a fully point-in-time,
free-data construction of PEAD and a same-industry lead--lag signal; (2) a Fama--MacBeth separation
of the lead--lag effect from own-firm PEAD, which we show is necessary because earnings-season
clustering confounds the naive estimator; (3) a utility-weighted gate that judges signals by
economic P\&L rather than accuracy; (4) a risk-parity combination and a regime-aware sizing layer;
and (5) an evaluation protocol centered on \emph{durability}. The entire pipeline is released as
open code.

\paragraph{Stance on honesty.} Following the discipline of the overfitting literature, we
pre-register falsification thresholds and report decay explicitly. Where the evidence does not
support a directional claim, we say so; a correctly-inferred null is reported as such.

\section{Data and point-in-time construction}
\label{sec:data}

We use SEC EDGAR XBRL company facts and submissions, which provide, for each disclosure, the
\emph{filing date}---rendering the data point-in-time by construction. Quarterly diluted earnings
per share are extracted with tag fallbacks; the fourth fiscal quarter is derived as the annual
figure less the first three quarters. Prices are adjusted closes for an \emph{as-traded} universe
(historical index membership) to avoid survivorship bias. All positions are formed with a strict
one-day execution lag ($T{+}1$) relative to the filing date.

\begin{definition}[Standardized unexpected earnings]\label{def:sue}
For firm $i$ and fiscal quarter $t$, with $E_{i,t}$ the reported quarterly EPS,
\[
\SUE_{i,t}=\frac{E_{i,t}-E_{i,t-4}}{\sigma_{i,t}},\qquad
\sigma_{i,t}=\operatorname{std}\bigl(\{E_{i,s}-E_{i,s-4}\}_{s=t-W}^{t-1}\bigr),
\]
where $E_{i,t-4}$ is the same quarter one year prior (seasonal random walk) and $W$ a trailing
window (default $W=8$). $\SUE$ requires no analyst data and is point-in-time.
\end{definition}

\section{Signals}
\label{sec:signals}

\subsection{Own-firm PEAD}
The own-firm signal is the cross-sectionally standardized $\SUE_{i,t}$, active over a holding
window $[\,T{+}1,\,T{+}1{+}H)$ after the announcement, with $H$ chosen from the measured decay
profile (Section~\ref{sec:eventstudy}).

\subsection{Same-industry lead--lag}
\begin{definition}[Peer-implied surprise]\label{def:peer}
Let $\mathcal{P}(i)$ be the same-SIC peers of firm $i$. The peer-implied surprise of firm $j$ on
day $d$ is
\[
s^{\mathrm{peer}}_{j,d}=\sum_{i:\,j\in\mathcal{P}(i)} z(\SUE_{i})\,\mathbf{1}\{d\in[\tau_i{+}1,\tau_i{+}1{+}H)\},
\]
where $\tau_i$ is firm $i$'s announcement day and $z(\cdot)$ the cross-sectional standardization.
\end{definition}
This operationalizes the economic-link predictability of Cohen and Frazzini (2008) at the industry
level (Menzly and Ozbas, 2010).

\section{Methodology}
\label{sec:method}

\subsection{Event-study drift profile and decay}
\label{sec:eventstudy}
\begin{definition}[Drift profile and cumulative drift]\label{def:drift}
With $z_e$ the standardized surprise of event $e$ and $\AR_{e,k}$ the market-adjusted abnormal
return $k$ days after entry,
\[
b(k)=\frac{\sum_e z_e\,\AR_{e,k}}{\sum_e z_e^2},\qquad B(h)=\sum_{k=1}^{h}b(k).
\]
\end{definition}
We fit $b(k)\approx b_0 e^{-k/\tau}$ on the leading window and report the half-life
$t_{1/2}=\tau\ln 2$ and the horizon $h^\star$ at which marginal drift loses significance.

\todo{Insert measured $b(k)$, $B(h)$, $t_{1/2}$, $h^\star$ on real data; Figure with the three
panels (drift, cumulative, durability).}

\subsection{Calendar-time inference}
Because holding windows overlap and announcements cluster in earnings season, event-time standard
errors are badly understated. We therefore form a daily calendar-time long--short portfolio and
apply Newey--West (1987) standard errors with lag equal to the holding horizon. The headline
statistics are the annualized information ratio $\IR=\sqrt{252}\,\bar r/\sigma_r$ and the
Newey--West $t$-statistic.

\subsection{Separating lead--lag from own-firm PEAD: Fama--MacBeth}
Naively measuring the peer effect double-counts own-firm PEAD because firm $j$ typically announces
near its peers. We instead run, each day $d$, the cross-sectional regression
\[
\AR_{\cdot,d}=\alpha_d+\beta^{\mathrm{own}}_d\,s^{\mathrm{own}}_{\cdot,d}
+\beta^{\mathrm{peer}}_d\,s^{\mathrm{peer}}_{\cdot,d}+\varepsilon_{\cdot,d},
\]
and treat the time series $\{\beta^{\mathrm{peer}}_d\}$ as the lead--lag factor return (Fama and
MacBeth, 1973), with Newey--West inference. The own coefficient absorbs PEAD, isolating the
genuine cross-firm effect.

\todo{Report $\IR$/$t_{\mathrm{NW}}$ of $\beta^{\mathrm{peer}}$ with vs.\ without the own control;
the gap quantifies the confound.}

\subsection{Utility-weighted gate}
\begin{definition}[Signal utility]\label{def:util}
For a signal $S$ with cross-sectional weights $w_{i,d}=z(S)_{i,d}/\sum_i|z(S)_{i,d}|$ and next-day
returns $r_{i,d+1}$, the daily utility is $u_d=\sum_i w_{i,d}\,r_{i,d+1}$ (a weight-times-return
objective in the spirit of Jane Street's Kaggle scoring).
\end{definition}
\begin{proposition}[Gate equals portfolio Sharpe]\label{prop:gate}
The annualized utility ratio $\sqrt{252}\,\bar u/\sigma_u$ equals the information ratio of the
signal's dollar-neutral long--short portfolio. Hence gating on utility is equivalent to gating on
risk-adjusted economic P\&L, not on classification accuracy.
\end{proposition}
\todo{One-line proof (definitional) + numerical cross-check on synthetic and real data.}
A signal passes if its out-of-sample $\IR>0$, $|t_{\mathrm{NW}}|\ge t^\star$ (default $t^\star=2$,
with the Deflated Sharpe Ratio and PBO reported alongside), and its late-subsample $\IR$ does not
collapse.

\subsection{Risk-parity combination}
\begin{definition}[Equal-risk-contribution weights]\label{def:erc}
Given surviving signals with portfolio return covariance $\Sigma$, the combination weights solve
$w_i\,(\Sigma w)_i=\text{const}$ for all $i$ (equal risk contribution). The inverse-volatility
special case is $w_i\propto 1/\sigma_i$.
\end{definition}
Weighting by risk---not capital---prevents a single noisy signal from dominating and adds no fitted
parameters, avoiding a fresh layer of overfitting.

\subsection{Regime-aware sizing}
\begin{definition}[Volatility-target multiplier]\label{def:regime}
With $\hat\sigma_d$ the trailing realized volatility of the (market or book) state and target
$\sigma^\ast$, the exposure multiplier is
$m_d=\min\!\bigl(\sigma^\ast/\hat\sigma_{d-1},\,c\bigr)$, capped at $c$ and lagged one day.
\end{definition}
This sizes the book down in high-volatility regimes and up in calm ones, in the spirit of Moreira
and Muir (2017); it controls risk and does not forecast direction.

\section{Evaluation protocol}
\label{sec:eval}
We use strictly walk-forward training/testing with an embargo; combinatorial purged
cross-validation and the probability of backtest overfitting (L\'opez de Prado, 2018); the Deflated
Sharpe Ratio (Bailey and L\'opez de Prado, 2014) to adjust for multiple trials; Newey--West
inference throughout; and year-by-year durability as the primary robustness diagnostic. Falsifi
cation thresholds are fixed in advance.

\section{Results}
\label{sec:results}
\todo{Tables: (1) per-signal gate table (IR, $t_{\mathrm{NW}}$, mean IC, early/late IR, DSR, PBO);
(2) lead--lag with/without own control; (3) combined book raw vs.\ regime-scaled (IR, vol, maxDD);
(4) durability by year. Figures: drift/decay, lead--lag cumulative + durability, engine equity +
exposure multiplier. Insert real-data numbers; synthetic validation figures may appear in an
appendix as estimator sanity checks, clearly labeled as such.}

\section{Scope and limitations}
\label{sec:scope}
Free daily data caps achievable edge; results are expected to be weak and possibly null, and we
report them as such. Time-series $\SUE$ avoids analyst-consensus dependence but omits the
revision channel. The same-industry map (SIC) is coarse relative to true supply-chain links; the
lead--lag effect is the more fragile of the two signals and most exposed to spurious structure.
Point-in-time fidelity depends on XBRL coverage (broadly complete only from about 2011).

\section{Conclusion}
\label{sec:conclusion}
The pipeline turns two classical, mechanism-backed anomalies into point-in-time, free-data signals,
separates a genuine cross-firm effect from its own-firm confound, combines weak signals by risk
parity, and sizes the result by regime---while making honest, durability-centered evaluation the
core deliverable. The value is reproducible methodology, whatever the sign of the final result.

\begin{thebibliography}{99}
\small
\bibitem{ballbrown1968} Ball, R., and Brown, P. (1968). An empirical evaluation of accounting income numbers. \textit{Journal of Accounting Research} 6(2), 159--178.
\bibitem{bt1989} Bernard, V. L., and Thomas, J. K. (1989). Post-earnings-announcement drift: delayed price response or risk premium? \textit{Journal of Accounting Research} 27, 1--36. \href{https://doi.org/10.2307/2491062}{doi:10.2307/2491062}.
\bibitem{fos1984} Foster, G., Olsen, C., and Shevlin, T. (1984). Earnings releases, anomalies, and the behavior of security returns. \textit{The Accounting Review} 59(4), 574--603.
\bibitem{cf2008} Cohen, L., and Frazzini, A. (2008). Economic links and predictable returns. \textit{The Journal of Finance} 63(4), 1977--2011. \href{https://doi.org/10.1111/j.1540-6261.2008.01379.x}{doi:10.1111/j.1540-6261.2008.01379.x}.
\bibitem{mo2010} Menzly, L., and Ozbas, O. (2010). Market segmentation and cross-predictability of returns. \textit{The Journal of Finance} 65(4), 1555--1580. \href{https://doi.org/10.1111/j.1540-6261.2010.01578.x}{doi:10.1111/j.1540-6261.2010.01578.x}.
\bibitem{mm2017} Moreira, A., and Muir, T. (2017). Volatility-managed portfolios. \textit{The Journal of Finance} 72(4), 1611--1644. \href{https://doi.org/10.1111/jofi.12513}{doi:10.1111/jofi.12513}.
\bibitem{hlz2016} Harvey, C. R., Liu, Y., and Zhu, H. (2016). $\ldots$ and the cross-section of expected returns. \textit{Review of Financial Studies} 29(1), 5--68. \href{https://doi.org/10.1093/rfs/hhv059}{doi:10.1093/rfs/hhv059}.
\bibitem{bldp2014} Bailey, D. H., and L\'opez de Prado, M. (2014). The Deflated Sharpe Ratio: correcting for selection bias, backtest overfitting, and non-normality. \textit{The Journal of Portfolio Management} 40(5), 94--107. \href{https://doi.org/10.3905/jpm.2014.40.5.094}{doi:10.3905/jpm.2014.40.5.094}.
\bibitem{ldp2018} L\'opez de Prado, M. (2018). \textit{Advances in Financial Machine Learning}. Wiley.
\bibitem{fm1973} Fama, E. F., and MacBeth, J. D. (1973). Risk, return, and equilibrium: empirical tests. \textit{Journal of Political Economy} 81(3), 607--636.
\bibitem{nw1987} Newey, W. K., and West, K. D. (1987). A simple, positive semi-definite, heteroskedasticity and autocorrelation consistent covariance matrix. \textit{Econometrica} 55(3), 703--708.
\bibitem{jt1993} Jegadeesh, N., and Titman, S. (1993). Returns to buying winners and selling losers: implications for stock market efficiency. \textit{The Journal of Finance} 48(1), 65--91. \href{https://doi.org/10.1111/j.1540-6261.1993.tb04702.x}{doi:10.1111/j.1540-6261.1993.tb04702.x}.
\bibitem{lm2006} Livnat, J., and Mendenhall, R. R. (2006). Comparing the post-earnings-announcement drift for surprises calculated from analyst and time-series forecasts. \textit{Journal of Accounting Research} 44(1), 177--205. \href{https://doi.org/10.1111/j.1475-679X.2006.00196.x}{doi:10.1111/j.1475-679X.2006.00196.x}.
\end{thebibliography}

\end{document}
